\section{Conclusion}
\label{sec:conclusion}

We introduced \PTA{}, a deliberately minimal two-stage framework that couples a fixed open-loop probing phase with a latent task-conditioning encoder and a belief-conditioned residual manipulation policy, and a reproducible cross-material edge-push benchmark in Genesis spanning sand, snow, and elastoplastic media (55pp of scripted-transfer spread). All methods are trained on sand only and evaluated zero-shot on the four other splits.

The empirical finding is conditional. On elastoplastic (recoverable rebound), \PTA{} delivers a 14.7pp transfer gain, a 14.7pp spill reduction, and a 16.7pp success-rate gain over a matched reactive PPO baseline, with matched ablations placing both probe and encoder \emph{below} the reactive baseline when either is removed. On the four other splits the same architecture underperforms by 13--22pp. We do not claim broad cross-material robustness. We offer the asymmetry as evidence consistent with a \emph{recoverable-deformation hypothesis}---probing appears beneficial when probe-induced perturbations relax on the task timescale and detrimental when they do not---and treat that hypothesis as a candidate boundary condition to be tested in future work, not as a demonstrated principle. A privileged-teacher baseline that collapses to 0\% on elastoplastic across three seeds is consistent with a complementary observation: in our sand-only-trained setup, an in-episode probing trace was a more useful conditioning signal than a passive parameter vector.

\PTA{} is therefore best read as a candidate tool for recoverable-response settings, scoped by an empirically observable boundary condition. Existing reactive policies remain appropriate for irreversibly perturbed (granular non-cohesive) media, and a stack handling mixed materials should include a coarse material-class router. Natural extensions include (i) testing the hypothesis on additional viscoelastic and granular materials, (ii) replacing the fixed probe with a learned information-gain probe, (iii) a gating mechanism that decides \emph{when} to probe, (iv) reusing the encoder across tasks, and (v) sim-to-real deployment on tactile-equipped hardware. Code, configurations, and seeds will be released upon acceptance.
